\scanpage{18}
affine, on constate que les ennuis proviennent du fait que le passage à la limite projective ne commute pas à la localisation par $p$
\diagram{\begin{tikzcd}[ampersand replacement=\&]
W \arrow[r] \& W[1/p]=K !
\end{tikzcd}}.

Je propose d'essayer l'approche suivante, consistant à remplacer le site cristallin par un site plus gros, formé avec des ``épaississements'' pas purement infinitésimaux : objets les triples
\[
(U_0,\unsure{\widetilde V},i:U_0\longrightarrow\unsure{\widetilde V}),
\]
où $U_0$ un ouvert affine de $X_0$, \inserted{(d'anneau $A_0$)}, $\unsure{\widetilde V}$ un ``\deleted{schéma}\inserted{espace} affine de W-M sur $W$'', \deleted{i.e. le} complété d'une \unsure{Alg.} de W-M $\underline B$ sur $W$,
\diagram{\begin{tikzcd}[ampersand replacement=\&,column sep=small]
U_0 \arrow[r,dash] \& \unsure{\widetilde V}
\end{tikzcd}}
une \deleted{immersion} morphisme \deleted{définie} induisant une \uline{immersion nilpotente}
\diagram{\begin{tikzcd}[ampersand replacement=\&]
U_0 \arrow[r] \& \widetilde V\times_S\unsure{s}
\end{tikzcd}}
(\unsure{alors} $i$ correspond à un $W$-hom surjectif
\diagram{\begin{tikzcd}[ampersand replacement=\&,column sep=small]
B \arrow[r] \& A_0
\end{tikzcd}}
tel que
\diagram{\begin{tikzcd}[ampersand replacement=\&,column sep=small]
B/pB \arrow[r,dash] \& A_0
\end{tikzcd}}
soit à noyau nilpotent).
Comme topologie, une variante convenable de la zariskienne (on peut se borner à \unsure{titre d'essai de ne regarder} que les traces des ouverts sur $X_0$).
\margin{17}
